% Abhakseite „Das kann ich“ – Zone nur im Gesamt (\mitzone)
%% TODO Ich-kann-Sätze: Platzhalter wie in den Aufgabendateien
\begin{abhakseite}
\ifdefined\mitzone
\abhakgruppe{Kennst du schon}
\abhak{1}{Vierfeldertafel füllen und lesen (Rand, Feld, bedingte Angabe unterscheiden)}
\abhak{2}{Baumdiagramm und Pfadregeln (Produkt entlang des Pfades, Summe über Pfade)}
\abhak{3}{Anteile und Prozentsätze umrechnen, Quotienten von Anteilen bilden und kürzen}
\abhak{4}{Brüche mit Termen vergleichen (Zähler konstant, Nenner ändert sich) und einfache Bruchungleichungen lösen}
\abhak{5}{Graphen lesen (Funktionswert zu Stelle, Stelle zu Funktionswert, Randwerte)}
\abhak{6}{Vierfeldertafel füllen und lesen (Rand, Feld, bedingte Angabe unterscheiden) – Fehler finden}
\abhak{7}{Vierfeldertafel füllen und lesen (Rand, Feld, bedingte Angabe unterscheiden) – selbst rechnen}
\fi
\abhakgruppe{Einheit 1 · Der Quotient}
\abhak{8}{Der Quotient}
\abhak{9}{Der Quotient – Fehler finden, Begründen, Anwendung}
\abhakgruppe{Einheit 2 · Bayes}
\abhak{10}{Bayes}
\abhak{11}{Bayes – Fehler finden, Begründen, Darstellungswechsel, Anwendung}
\abhakgruppe{Einheit 3 · Mit Parameter}
\abhak{12}{Mit Parameter}
\abhak{13}{Mit Parameter – Fehler finden, Begründen, Anwendung}
\end{abhakseite}
